% !TeX root = GeoDynamics.tex
\section{Drift calculations}

Drift calculation is an accumulation of momentum changes during a simulation. The drift is:


\begin{equation}
\begin{aligned}
\mathbf{p}_{CurrentTotalMomentum} = \mathbf{p}_{ParticleVelocityMomentum} \\ + \mathbf{p}_{ ParticleSideInternalMomentum} + \mathbf{p}_{WallGhostInternalMomentum} + \mathbf{p}_{WallVelocityMomentumReservoir}
\end{aligned}
\end{equation}

Where velocity momentum, $\mathbf{p}_{ParticleVelocityMomentum} $, is the sum of all the particles translational momentum. Internal momentum, $\mathbf{p}_{ParticleSideInternalMomentum}$. Wall momentum \\ $\mathbf{p}_{ WallGhostInternalMomentum}$ is the value of the wall internal momentum. And \\ $\mathbf{p}_{WallVelocityMomentumReservoir}$ is the pseudo wall reaction velocity.

The total starting momentum $\mathbf{p}_{InitialTotalMomentum}$, minus the current total momentum \\  $\mathbf{p}_{CurrentTotalMomentum}$ is drift $\mathbf{p}_{Drift}$.

Reporting variable perform a running sum where as each particles momentum values are recorded by adding them to previous reporting values. This means for each particle, at a time step, the kinetic momentum is added to the global variable $\mathbf{p}_{ ParticleVelocity Momentum}$, and likewise for the other variables. They are independently assigned to reporting variables. At the end of the time step the filled reporting variables are used to calculate the drift.  

Drift tells us how much the total system is varying from the current system and shows us the accumulation/system errors. The  accumulation/system errors become very important in large scale simulation of millions of particles.

The model of wall internal momentum may not match the actual dynamics but it is viewed as follows: The wall/ghost internal momentum ,$\mathbf{p}_{WallGhostInternalMomentum}$ is equal and opposite to the particle-side,$\mathbf{p}_{ParticleSideInternalMomentum}$  wall-contact internal momentum. 

The model of particle internal momentum may not match the actual dynamics but it is viewed as follows:Each particles internal momentum accumulates during compression and feeds back during rebound. During compression it is a vector opposite the velocity vector. In this way the internal momentum is fed back in the same direction as rebound. 

Particle-particle collisions result in four offsetting vectors, two opposing velocity momentum vectors and two opposing internal momentum vectors. Particle-wall collisions differ mathematically since the wall does not move. There are then only three vectors - two internal momentum vectors and only one velocity momentum vector. This leaves an imbalance of one velocity momentum vector. To account for this we create a wall-side velocity momentum reservoir $\mathbf{p}_{WallVelocityMomentumReservoir}$.


